\documentclass[12pt,letterpaper]{article}



%%
%% Required packages:
%%
\usepackage{etex}
\usepackage{amsmath}
\usepackage{amssymb}
\usepackage{amstext}
\usepackage{amsthm}
\usepackage{fancyhdr,fancybox}
\usepackage{mathrsfs} % need for \mathscr command
\usepackage{multicol}
\usepackage[normalem]{ulem}
%\usepackage{pictex,pictexplus}
% \usepackage{pictexwd}
\usepackage{rotating}
\usepackage{url}
\usepackage{xfrac}
\usepackage{booktabs}
\usepackage{color,soul}
\usepackage{comment}

\usepackage{hyperref}
\hypersetup{
    colorlinks=true,     % false: boxed links; true: colored links
    linkcolor=blue,      % color of internal links
    citecolor=blue,      % color of links to bibliography
    filecolor=blue,      % color of file links
    urlcolor=blue        % color of external links
}


\usepackage{tikz}
\usetikzlibrary{backgrounds,calc}

%%
%% Common formatting commands
%%
\setlength{\parindent}{0in}
\setlength{\textwidth}{7in}
\setlength{\evensidemargin}{-0.25in}
\setlength{\oddsidemargin}{-0.25in}
\setlength{\parskip}{.5\baselineskip}
\setlength{\topmargin}{-0.5in}
\setlength{\textheight}{9in}

%               Problem and Part
\newcounter{problemnumber}
\newcounter{partnumber}
\newcounter{subpartnumber}
\newcommand{\Problem}{%
  \refstepcounter{problemnumber}%
  \setcounter{partnumber}{0}%
  \setcounter{subpartnumber}{0}%
  \item[\makebox{\hfil\textbf{\theproblemnumber.}\hfil}]%
  }
\newcommand{\InProblem}[2][1.6]{%
  \refstepcounter{problemnumber}%
  \setcounter{partnumber}{0}%
  \setcounter{subpartnumber}{0}%
  \textbf{\theproblemnumber.}\ \ \parbox[t]{#1in}{#2}%
}
\newcommand{\InSmallProblem}[1]{\InProblem[1.05]{#1}}
\newcommand{\NoProblem}[1]{%
  \mbox{\hfil\textbf{#1}\hfil}%
  }
\newcommand{\UnProblem}{%
  \refstepcounter{problemnumber}%
  \setcounter{partnumber}{0}%
  \setcounter{subpartnumber}{0}%
  \mbox{\hfil\textbf{\theproblemnumber.}\hfil}%
  }
\newcommand{\InFlexProblem}[1]{%
  \refstepcounter{problemnumber}%
  \setcounter{partnumber}{0}%
  \setcounter{subpartnumber}{0}%
  \textbf{\theproblemnumber.}\ \ #1%
}
%%  Part:
\newcommand{\Part}{%
  \renewcommand{\thepartnumber}{(\alph{partnumber})}%
  \refstepcounter{partnumber}
  \setcounter{subpartnumber}{0}%
  \item[(\alph{partnumber})]%
}
\newcommand{\InPart}[2][1.75]{%
  \renewcommand{\thepartnumber}{(\alph{partnumber})}%
  \refstepcounter{partnumber}%
  \setcounter{subpartnumber}{0}%
  (\alph{partnumber})\ \ \parbox[t]{#1in}{#2}%
}
\newcommand{\UnPart}{%
  \refstepcounter{partnumber}%
  \renewcommand{\thepartnumber}{(\alph{partnumber})}%
  \setcounter{subpartnumber}{0}%
  (\alph{partnumber})%
}
\newcommand{\InLargePart}[1]{\InPart[3.05]{#1}}
\newcommand{\InFullPart}[1]{\InPart[6.2]{#1}}
\newcommand{\InSmallPart}[1]{\InPart[1.05]{#1}}
%% SubPart:
\newcommand{\SubPart}{%
  \refstepcounter{subpartnumber}%
  \item[(\roman{subpartnumber})]%
  }
\newcommand{\InSubPart}[2][2.25]{%
  \stepcounter{subpartnumber}%
  (\roman{subpartnumber})\ \ \parbox[t]{#1in}{#2}%
}
\newcommand{\InSmallSubPart}[1]{\InSubPart[1.5]{#1}}
\newcommand{\InTinySubPart}[1]{%
  \stepcounter{subpartnumber}%
  (\roman{subpartnumber})\ \ \makebox{#1}%
}
\newcommand{\InMySubPart}[2][2.25]{%
  \stepcounter{subpartnumber}%
  \parbox[t]{#1in}{\makebox[0.3in][l]{(\roman{subpartnumber})}\ \ {#2}}%
}
\newcommand{\TrueFalse}{\textbf{T}\ \ \textbf{F}\ \ }
\newcommand{\TFPart}[1]{% 
  \stepcounter{partnumber}%
  \setcounter{subpartnumber}{0}%
  \item[(\alph{partnumber})] \TrueFalse\parbox[t]{5.5in}{#1}}

\newcommand{\Points}[1]{(#1 points) \ }
\newcommand{\Point}[1]{(#1 point) \ }


%% For Answers:
\newcounter{answernumber}
\newcommand{\Answer}{%
  \refstepcounter{answernumber}%
  \setcounter{partnumber}{0}%
  \setcounter{subpartnumber}{0}%
  \item[\makebox{\hfil\textbf{\theanswernumber.}\hfil}]%
  }


% Setting up the headers for the first page
%\chead{} % will default to what is typed in the file.
\fancyhead[L]{\textsc{Math 22a}}
\fancyhead[R]{\textsc{Fall 2024}}
\fancyfoot[R]{
	\vspace*{\fill}\footnotesize\textcopyright{} \emph{Harvard Math 22a}	
}
\renewcommand{\headrulewidth}{0pt}

%\fancyhead[C]{}
%	\chead{}	
%	\lhead{}
%	\rhead{}
%	\cfoot{}


% Putting copyright on the footer of each page
%\fancypagestyle{secondpage}{
%	\fancyhead[C]{TESTing again} % change this to be blank
%		%\chead{}	
%	\fancyhead[L]{bears} % change this to be blank
%	\fancyhead[R]{dogs} % change this to be blank
%	\fancyfoot[R]{ %keeps the footer the same
%		\vspace*{\fill}\footnotesize\textcopyright{} \emph{Harvard Math 22a}	
%	}
%}
%\renewcommand{\headrulewidth}{0pt}
%\pagestyle{fancy}
\pagestyle{empty}


%\lhead{}
%\rhead{}
%\rfoot{\vspace*{\fill}\footnotesize\textcopyright{} \emph{Harvard Math 22a}}
%\renewcommand{\headrulewidth}{0pt}
%\pagestyle{fancy}




%%
%% Common shortcut commands:
%%
\newcommand{\ds}{\displaystyle}
\newcommand{\sgn}{\operatorname{sgn}}
\renewcommand{\Pr}{\operatorname{P}}
\newcommand{\CPr}[3][{}]{\Pr_{\text{#1}}\left(#2\,|\,#3\right)}

\newcommand{\C}{\mathbb{C}}
\newcommand{\R}{\mathbb{R}}
\newcommand{\Z}{\mathbb{Z}}
\newcommand{\N}{\mathbb{N}}
\newcommand{\Q}{\mathbb{Q}}
\newcommand{\D}{\mathbb{D}}

\newcommand{\rref}{\operatorname{rref}}
\newcommand{\im}{\operatorname{im}}
\newcommand{\rank}{\operatorname{rank}}
\newcommand{\diag}{\operatorname{diag}}
\newcommand{\Span}{\operatorname{span}}
%\renewcommand{\vec}[1]{\mathbf{#1}}
\newcommand{\charpoly}{\mathscr{P}}
\newcommand{\nicefrac}[2]{\sfrac{#1}{#2}} % using xfrac package
\newcommand{\topology}{\mathcal{T}}
\newcommand{\Inn}{\operatorname{Inn}}

\newcommand{\blankbox}[2]{\fbox{\rule{#1}{0in}\rule{0in}{#2}}}

\newenvironment{myindentpar}[1]%
 {\begin{list}{}%
         {\setlength{\leftmargin}{#1}
          \setlength{\rightmargin}{\leftmargin}}%
         \item[]%
 }
 {\end{list}}
\newtheorem{theorem}{Theorem}
\newtheorem*{NStheorem}{Theorem}
\newtheorem{cor}[theorem]{Corollary}
\newtheorem*{claim}{Claim}
\newtheorem{defn}{Definition}

\newenvironment{amatrix}[1]{%
  \left(\begin{array}{@{}*{#1}{c}|c@{}}
}{%
  \end{array}\right)
}

%--------grstep
% For denoting a Gauss' reduction step.
% Use as: \grstep{\rho_1+\rho_3} or \grstep[2\rho_5 \\ 3\rho_6]{\rho_1+\rho_3}
\newcommand{\grstep}[2][\relax]{%
   \ensuremath{\mathrel{
       {\mathop{\longrightarrow}\limits^{#2\mathstrut}_{
                                     \begin{subarray}{l} #1 \end{subarray}}}}}}
\newcommand{\swap}{\leftrightarrow}




\chead{\Large\textbf{Linear and Matrix Transformations, $\vec\S$1.9}}% 

\begin{document}
\thispagestyle{fancy}

Recall from last class:

\begin{itemize}

\item Define $T : \R^n \to \R^m$ by $T(\vec{x})=A\vec{x}$ where $A$ is an $m$ by $n$ matrix. Recall we call $\R^n$ the domain and $\R^m$ the codmain. The set of outputs, $\{T(\vec{x}):\vec{x}\in \R^n\}$, is the image or the range of $T$.
\item $T : \R^n \to \R^m$ is {\bf linear} if 
\begin{enumerate}
	\item $T(\vec{u}+\vec{v})=T(\vec{u})+T(\vec{v})$ for all $\vec{u},\vec{v}$ in $\R^n$,
	\item $T(c\vec{u})=cT(\vec{u})$ for all $\vec{u}$ in $\R^n$ and $c$ in $\R$.
\end{enumerate}
\item {\bf Properties of Linear Transformations.}
\begin{enumerate}
\item If $T$ is linear, then $T(\vec{0})=\vec{0}$.
\item If $T$ is linear, then $T(c\vec{u}+d\vec{v})=cT(\vec{u})+dT(\vec{v})$.
\end{enumerate} 
\item All matrix transformations are linear.
\item If $T:\mathbb{R}^n \to \mathbb{R}^m$ is given by $T(\vec{x})=A\vec{x}$, then the image of $T$ is the span of the columns of $A$; namely, $$\text{Im}(T)= \text{Span}\{ \vec{a}_1, \dots, \vec{a}_n\}.$$
\item We also saw various geometric transformations, including projections, dilations, and sheers.
\end{itemize}

\newpage

\begin{enumerate}

\Problem Suppose $T: \R^2 \to \R^3$ is linear and $$T\left( \begin{bmatrix} 1\\ 0\\ \end{bmatrix} \right)=\begin{bmatrix} 5\\ -7\\ 2\\ \end{bmatrix}, \text{ and } T\left( \begin{bmatrix} 0\\ 1\\ \end{bmatrix} \right)=\begin{bmatrix} -3\\ 8\\ 0\\ \end{bmatrix}.$$ What is $$T\left( \begin{bmatrix} a\\ b\\ \end{bmatrix} \right)?$$

\vfill

%\Problem Suppose $T: \R^2 \to \R^3$ is linear and $

{\bf Theorem 10:} Let $T: \R^n \to \R^m$ be linear. Then there exists a unique matrix $A$ such that $T(\vec{x})=A\vec{x}$ for all $\vec{x} \in \R^n$.

\Problem Prove the Theorem.

\vfill

\vfill 

\newpage

\Problem Let $T: \R^2 \to \R^2$ be the linear transformation given by rotating $\vec{x}$ by $\theta$ counterclockwise. Find the associated matrix of $T$.

\vfill

\Problem Let $T: \R^4 \to \R^3$ be defined by $T(\vec{x})=A\vec{x}$ where $$A=\begin{bmatrix}
1 & -4 & 8 &1\\
0 & 2 & -1 &3\\
0 & 0 &0 & 5\\
\end{bmatrix}.$$ Is $T$ injective and/or surjective?

\vfill

\newpage

{\bf Theorem 11:} Let $T: \R^n \to \R^m$ be linear. Then $T$ is injective if and only if $T(\vec{x})=\vec{0}$ has only the trivial solution.

\Problem Prove the theorem.

\vfill 

{\bf Theorem 12:} Let $T: \R^n \to \R^m$ be linear, and let $A$ be the associated matrix. Then 
\begin{enumerate}
\item $T$ is surjective if and only if the columns of $A$ span $\R^m$.
\item $T$ is injective if and only if the columns of $A$ are linearly independent. 
\end{enumerate}

\Problem Prove the theorem.

\vfill 

\end{enumerate}

\end{document}

\newpage
\chead{\Large\textbf{Linear and Matrix Transformations -- Answers / Solutions}}%
\lhead{}
\rhead{}
\thispagestyle{fancy}

\begin{enumerate}
  \Answer Using linearity, we get 
  
  \begin{align*}
  T\left( \begin{bmatrix} a\\ b\\ \end{bmatrix} \right) &= T\left( a \begin{bmatrix} 1\\ 0\\ \end{bmatrix}+ b\begin{bmatrix} 0\\1\\ \end{bmatrix} \right)\\
  & = a T\left( \begin{bmatrix} 1\\ 0\\ \end{bmatrix} \right)+ bT\left( \begin{bmatrix} 5\\ -7\\ 2\\ \end{bmatrix} \right)\\
  &= a \begin{bmatrix} 1\\ 0\\ \end{bmatrix} + b\begin{bmatrix} -3\\ 8\\0 \end{bmatrix}\\
  &= \begin{bmatrix}
  	5 & -3\\
  	-7 & 8\\
  	2 &0\\
  	\end{bmatrix} \begin{bmatrix} a\\ b\\ \end{bmatrix}.
  \end{align*}

Thus $T$ is also a matrix transformation. A natural question is whether every linear transformation is a matrix transformation.
  
  \Answer   We proceed as in the last example. First Let $\vec{e}_j$ be the vector in $\R^n$ with 1 in the $j$-th position and 0 elsewhere. Let $x \in \R^n$, then $$\vec{x}=x_1 \vec{e}_1+\dots+x_n \vec{e}_n.$$ Thus using the linearity of $T$, we get
  \begin{align*}
  T(\vec{x}) & = x_1 T(\vec{e}_1)+\dots+x_n T(\vec{e}_n)\\
  &= \left[T(\vec{e}_1) \; \dots \; T(\vec{e}_n)\right]\vec{x}.
  \end{align*}
   Thus $T(\vec{x})=A\vec{x}$ where $A=\left[T(\vec{e}_1) \; \dots \; T(\vec{e}_n)\right]$.
   
   In order to prove uniqueness, we suppose there is another matrix $B$ such that $T(\vec{x})=B\vec{x}$ for all $x\in \R^n$. We want to show that $A=B$. Let $B=[\vec{b}_1 \; \dots \; \vec{b}_n]$. Since $A\vec{e}_j=B\vec{e}_j$, we get $\vec{b}_j=T(\vec{e}_j)$. Thus $A$ and $B$ have the same columns, and thus $A=B$ and $A$ is unique.
  \Answer Rotating $\vec{e}_1$ and $\vec{e}_2$ by $\theta$ gives the associated matrix of $T$ is 
  $\begin{bmatrix}
  \cos{\theta} & -\sin{\theta}\\
  \sin{\theta} & \cos{\theta}\\
  \end{bmatrix}$.
  
  \Answer We have a pivot position in each row, so $T$ is onto (surjective). $A\vec{x}=\vec{0}$ has free variables, so $T$ is not one-to-one (injective).
  
  \Answer Let $T:\R^n \to \R^m$ be linear.
  
  $\implies$ Suppose $T$ is one-to-one (injective). Since $T$ is linear, we know $\vec{x}=\vec{0}$ is a solution. Since$T$ is injective, we know $T(\vec{x})=\vec{0}$ has at most one solution. Thus $T(\vec{x})=\vec{0}$ has only the trivial solution.
  
  $\impliedby$ Suppose $T(\vec{x})=\vec{0}$ has only the trivial solution and $T(\vec{u})=T(\vec{v})$. Since $T$ is linear, we get $T(\vec{u}-\vec{v})=\vec{0}$. Since $T(\vec{x})=\vec{0}$ has only the trivial solution, we know that $\vec{u}-\vec{v}=\vec{0}$, and hence $\vec{u}=\vec{v}$. Therefore $T$ is one-to-one or injective.
  
  \Answer Let $T: \R^n \to \R^m$ be linear and let $A$ be the associated matrix.
  
  {\bf Proof of (a):} $T$ is surjective means that for all $\vec{b}\in \R^m$ there exists $\vec{x}\in \R^n$, or equivalently $A\vec{x}=\vec{b}$ is consistent. By Theorem 4, $A\vec{x}=\vec{b}$ is consistent for all $\vec{b}$ if and only if the columns of $A$ span $\R^m$.
  
    {\bf Proof of (b):} By Theorem 11, $T$ is one-to-one if and only if $T(\vec{x})=\vec{0}$ has only the trivial solution. Thus the columns of $A$ must be linearly independent.
  
\end{enumerate}

\end{document}


%%% Local ables: 
%%% mode: latex
%%% TeX-master: t
%%% End: 
